\chapter{Hemorrhoid Treatment}
\index{Hemorrhoid Treatment}

\section*{Definition and Overview}
Hemorrhoids (also spelled haemorrhoids) are swollen, inflamed veins in the lower rectum and anus that cause bleeding, itching, and discomfort. Hemorrhoid treatment ranges from simple self-care measures such as fibre, fluids, and over-the-counter creams for mild symptoms, to office procedures such as banding and injection for persistent internal hemorrhoids, and surgery for the most severe cases. Because most hemorrhoids are mild and related to straining, diet and bowel habits are the foundation of both treatment and prevention. Treatment is chosen according to the type, severity, and persistence of symptoms.

\section*{Epidemiology}
Hemorrhoids are very common, affecting up to half of adults at some point. They become more prevalent with age, peaking between 45 and 65 years, and are equally common in men and women, though pregnancy commonly triggers or worsens them. Most people with hemorrhoids have mild symptoms that respond to self-care, and only a minority need procedural treatment. Hemorrhoid surgery is one of the most commonly performed procedures, but is reserved for patients with severe, persistent, or thrombosed hemorrhoids.

\section*{Aetiology and Pathophysiology}
\begin{itemize}
    \item \textbf{Increased pressure:} straining at stool, prolonged sitting on the toilet, and constipation raise pressure in the anal veins
    \item \textbf{Pregnancy:} pressure from the growing uterus and hormonal changes enlarge anal veins
    \item \textbf{Obesity:} increased abdominal pressure contributes
    \item \textbf{Chronic diarrhoea:} frequent straining and irritation
    \item \textbf{Pathophysiology:} the anal cushions, normal vascular tissue in the anal canal, engorge and prolapse, forming internal or external hemorrhoids
    \item \textbf{Internal vs external:} internal hemorrhoids arise above the dentate line and usually bleed painlessly; external ones arise below it and are painful when thrombosed
    \item \textbf{Treatment targets:} relieve symptoms, reduce blood flow to the hemorrhoid, fix it in place, or remove it
\end{itemize}

\section*{Clinical Features}
\begin{itemize}
    \item Bright red bleeding on the toilet paper or in the bowl after a bowel motion
    \item Itching, irritation, and moisture around the anus
    \item A lump at the anus, either a prolapsed internal hemorrhoid or an external one
    \item Discomfort, aching, or a feeling of incomplete emptying
    \item Sudden severe pain with a tender, hard lump indicates a thrombosed external hemorrhoid
    \item Prolapsed hemorrhoids may need to be pushed back or may remain outside
\end{itemize}

\section*{Differential Diagnosis}
\begin{itemize}
    \item Anal fissures cause sharp pain and bleeding with bowel motions
    \item Perianal abscess and anal fistula cause pain, swelling, and discharge
    \item Colorectal cancer and polyps can cause rectal bleeding and must be considered
    \item Inflammatory bowel disease causes bleeding and diarrhoea
    \item Skin tags and other anal lumps are benign but distinct from hemorrhoids
    \item Any rectal bleeding, especially in older people or with changed bowel habit, warrants assessment
\end{itemize}

\section*{When to Seek Medical Care}
See a doctor for any rectal bleeding, a persistent anal lump, or symptoms that do not settle with self-care. New bleeding or a change in bowel habit in anyone over 50, or in anyone with a family history of bowel cancer, requires prompt investigation.

\textbf{Contact your local emergency services for:}
\begin{itemize}
    \item Heavy or persistent rectal bleeding with faintness or dizziness
    \item Passing large amounts of blood or black, tarry stools
    \item Severe anal pain with fever and swelling, suggesting abscess
    \item Inability to pass stool with abdominal pain and vomiting
\end{itemize}

\section*{Diagnosis and Investigations}
\begin{itemize}
    \item \textbf{Examination:} inspection of the anus and digital rectal examination
    \item \textbf{Proctoscopy and anoscopy:} direct visualisation of internal hemorrhoids
    \item \textbf{Colonoscopy:} recommended for rectal bleeding, especially in people over 50 or with risk factors, to exclude more serious causes
    \item \textbf{Assessment:} the type (internal or external), grade, and severity guide treatment
    \item \textbf{Follow-up:} surveillance according to findings
\end{itemize}

\section*{Management}
\begin{itemize}
    \item \textbf{Self-care:} increased fibre and fluids, softer stools, no straining, warm baths, and topical creams or suppositories
    \item \textbf{Fibre supplements and laxatives:} reduce straining and keep stools soft
    \item \textbf{Rubber band ligation:} a band is placed at the base of an internal hemorrhoid, which shrinks and falls off within days
    \item \textbf{Sclerotherapy:} injection of a chemical to shrink the hemorrhoid
    \item \textbf{Infrared coagulation and cautery:} heat-based office treatments
    \item \textbf{Thrombosed external hemorrhoid:} pain relief and, if very painful and early, surgical drainage of the clot
    \item \textbf{Surgery (hemorrhoidectomy):} excision of severe, prolapsing, or recurrent hemorrhoids
    \item \textbf{Stapled and laser techniques:} alternative surgical approaches with different trade-offs
\end{itemize}

\section*{Complications}
\begin{itemize}
    \item Thrombosis of an external hemorrhoid, causing sudden severe pain
    \item Prolapse and incarceration of internal hemorrhoids
    \item Chronic bleeding leading to iron-deficiency anaemia
    \item Infection and abscess formation, rarely
    \item Postoperative pain, bleeding, and difficulty passing urine after surgery
    \item Recurrence after treatment, especially without dietary change
    \item Delayed diagnosis of more serious causes of rectal bleeding
\end{itemize}

\section*{Prognosis and Outcomes}
Most hemorrhoids respond to dietary and lifestyle measures, and the majority of people never need more than self-care. Office procedures such as banding are effective for most internal hemorrhoids, with high success rates and minimal downtime. Hemorrhoidectomy is effective for severe disease, though recovery is painful for a week or two and recurrence is possible if straining continues. The long-term outlook is excellent when the underlying causes, particularly constipation and straining, are corrected.

\section*{Prevention and Self-Care}
\begin{itemize}
    \item Eat a high-fibre diet with plenty of fruit, vegetables, and wholegrains
    \item Drink plenty of water daily
    \item Do not strain at stool or sit on the toilet for long periods
    \item Respond promptly to the urge to open the bowels
    \item Exercise regularly and maintain a healthy weight
    \item Take warm baths to soothe symptoms
    \item Use soft toilet paper or wet wipes and avoid harsh rubbing
    \item Seek medical advice for persistent or bleeding hemorrhoids
\end{itemize}

\section*{Sources and Further Reading}
The following reputable sources provide further information for healthcare students and patients seeking reliable, current advice about hemorrhoid treatment.

\begin{itemize}
    \item Healthdirect. \textit{Haemorrhoids: treatment options.} https://www.healthdirect.gov.au/haemorrhoids
    \item The Gut Foundation. \textit{Haemorrhoids information.} https://gutfoundation.org.au/
    \item Bowel Cancer Australia. \textit{Rectal bleeding: don't assume it's hemorrhoids.} https://www.bowelcanceraustralia.org/
    \item NHS. \textit{Piles (haemorrhoids): treatment.} https://www.nhs.uk/conditions/piles/
    \item Mayo Clinic. \textit{Hemorrhoids: diagnosis and treatment.} https://www.mayoclinic.org/
    \item MSD Manual. \textit{Hemorrhoids.} https://www.msdmanuals.com/
\end{itemize}
